\documentclass[11pt]{article}
\usepackage{icmlko}
\begin{document}

\title{라벨 다섯 개면 버린 축을 되찾는다 --- 영샷이 못 하는 일을
몇 샷이 공짜로 한다}
\author{월드모델 실험실}
\date{노트 443 · 2026년 8월 2일}
\maketitle

\begin{abstract}
사흘 사이에 축 하나를 놓고 세 번 물었다. 노트 441 은
\texttt{entry\_friction} 이 \textbf{수집마다 부호가 뒤집히는} 축이라 안
배운 수집에서는 \textbf{버리는} 게 낫다는 걸 보였다(집 밖 두 짝에서
$+0.3174$ · $+0.4188$). 노트 442 는 \textbf{계산되는 표지로 되찾기}가 오히려
더 나쁘다는 걸 보였다. 남은 길은 하나였다 --- \textbf{새 수집의 라벨 몇
개로 부호를 직접 재는 것.} 영샷이 아니라 몇 샷이다.

\textbf{다섯 개면 된다.}

\begin{center}
\begin{tabular}{lrrr}
\hline
 & KR 만화 & 비게임 앱 \\
\hline
그대로 (축을 믿는다) & $+0.3206$ & $-0.0501$ \\
가림 --- 노트 441 이 채택 ($k{=}0$) & $+0.5336$ & $+0.2635$ \\
\hline
몇 샷 $k{=}5$ & $\mathbf{+0.5690}$ & $\mathbf{+0.4307}$ \\
몇 샷 $k{=}10$ & $+0.5897$ & $+0.4636$ \\
몇 샷 $k{=}40$ & $\mathbf{+0.5999}$ & $\mathbf{+0.4817}$ \\
\hline
오라클 부호 (전 행) & $+0.6000$ & $+0.4817$ \\
\hline
\end{tabular}
\end{center}

\textbf{라벨 다섯이면 버리는 것보다 낫고, 마흔이면 오라클과 소수점 넷째
자리까지 같다.} 두 짝 모두. 채점은 부호를 재는 데 쓴 $k$ 행을 \emph{뺀}
나머지에서 했다.
\end{abstract}

\begin{figure}[t]
\centering
\includegraphics[width=\textwidth]{figs/f1.pdf}
\caption{$k$ 를 늘릴 때 세 규칙. 굵은 파랑이 채택한 규칙 B 다. 초록 점선이
오라클, 빨강 점선이 노트 441 이 채택한 가림($k{=}0$), 회색 점선이 축을
그대로 믿는 것. \textbf{규칙 B 는 $k{=}5$ 에서 이미 빨강 위에 있고
$k{=}40$ 에서 초록에 닿는다.} 규칙 A 가 작은 $k$ 에서 처지는 것이 이
노트가 고친 실수다.}
\end{figure}

\section{되돌아갈 자리를 잘못 골랐었다}

처음 짠 규칙은 \emph{``부호가 음수면 뒤집고 아니면 \textbf{그대로 둔다}''}
였다(규칙 A). $k$ 가 작으면 부호 추정이 틀리는데, 틀렸을 때 떨어지는
자리가 \textbf{축을 그대로 믿는 자리}다 --- KR 에서 $+0.32$, 앱에서
$-0.05$. 그래서 $k{=}5$ 에서 규칙 A 는 $+0.4695$ 로 \textbf{가림보다도
못했다}.

고친 규칙 B 는 \emph{``음수면 뒤집고 아니면 \textbf{가린다}''} 이다. 되돌아갈
자리를 노트 441 이 이미 증명한 안전한 자리로 바꾼 것뿐인데 $k{=}5$ 에서
$+0.4695 \to \mathbf{+0.5690}$ 이다. \textbf{몇 샷 규칙의 값은 맞혔을 때가
아니라 틀렸을 때 어디로 떨어지느냐가 정한다.}

신뢰구간으로 문을 하나 더 단 규칙 C(\,$|$부호$| <$ SE 면 가린다\,)는 규칙 B
보다 모든 $k$ 에서 낮거나 같았다. 부호가 약할 때도 \emph{뒤집는 쪽}이
이득이다 --- 문을 달 이유가 없다.

\section{어디에 넣었나}

\texttt{F18\_bagboost.calibrate\_ooc(d, A, M, y)} 로 넣었다. \textbf{안
부르면 예전 그대로 가린다} --- 영샷 기본값은 안 바뀌고, 판은 이 변경도
원리상 못 본다(판이 채점하는 도메인은 전부 학습에 들어 있다). 가드 여섯을
붙였다: 집안 예보 최대 차 $0.00\mathrm{e}{+}00$ · 입력 불변 · 보정하면
달라지고 · 부호가 양수로 나오면 가림과 \emph{정확히} 같아진다.

\section{무엇을 산 것인가}

이 셋을 붙이면 이야기가 닫힌다.

\begin{itemize}
\item 축은 \textbf{정보가 없어서} 못 쓴 게 아니다 --- 도메인 안에서 값이
크다(모바일 $+0.648$).
\item 못 쓴 이유는 \textbf{어느 쪽으로 읽을지 몰라서}다. 집 안에서는
원핫이 알려 주고 \textbf{집 밖에는 그게 없다}.
\item 그 모르는 것은 \textbf{계산으로는 못 메우고}(노트 442) \textbf{라벨
다섯 개로 메워진다}.
\end{itemize}

새 IP 수집을 열 때의 실무 규칙이 여기서 나온다. \textbf{마흔 개만 라벨을
달면 뒤집히는 축이 제자리로 온다.} 그 뒤로는 더 달아도 안 는다.

\section*{남은 것}

\begin{itemize}
\item \textbf{다른 축도 뒤집히나.} 지금 뒤집기를 거는 축은
\texttt{entry\_friction} 하나다(다섯 중 부호가 안 맞는 유일한 축). 축을
늘리면 --- 36 축 중 표시자가 있는 것들 --- 몇 샷의 값이 더 커지나.
\item \textbf{k 를 축마다 나눠 쓰기.} 라벨 마흔은 축 하나에 쓰든 넷에
쓰든 마흔이다. 어디에 쓰는 게 나은지 안 재 봤다.
\item \textbf{CN 이 왜 안 오르나} --- 독을 안 먹었는데 $+0.2249$. 라벨
분포가 마지막 후보다.
\end{itemize}

\end{document}
